\section{Conclusões e trabalho futuro}
\label{sec:conclusoes}

\subsection{Conclusões}

Este trabalho apresentou uma revisão sistemática da literatura sobre mecanismos de
segurança para comunicações IoT baseadas em MQTT, conduzida segundo um protocolo
explícito --- perguntas de investigação, PICOC, \emph{strings} de pesquisa por base
de dados, critérios de inclusão, exclusão e qualidade, leitura em três passagens e
formulário de extração --- e relatada com o diagrama PRISMA. Foram analisados
\nIncluidos{} estudos primários (2014--2026) e \nSecundarios{} revisões
relacionadas.

As principais conclusões são as seguintes:

\begin{itemize}[leftmargin=*]
  \item O MQTT delega a segurança na implementação, e o TLS, por terminar no
    \emph{broker}, não protege as mensagens de um \emph{broker} comprometido nem de
    clientes com acesso legítimo, como mostram os estudos de análise de ameaças.
  \item A literatura organiza-se em oito categorias. A proteção leve fim-a-fim da
    mensagem é a mais numerosa, e a confidencialidade, a autenticação e a
    autorização são as propriedades mais tratadas.
  \item A frescura e a anti-repetição, a ligação da mensagem ao tópico, a política
    contra a redução do nível de proteção, a restrição chave--tópico e o não repúdio
    são pouco ou nada reportados, e a maioria das soluções continua a confiar no
    \emph{broker}.
  \item As avaliações medem o custo dos mecanismos, mas raramente o relacionam com
    os ataques efetivamente impedidos.
\end{itemize}

Com base nestas lacunas, propôs-se uma arquitetura de segurança em camadas
independentes e comutáveis, que combina defesa no transporte, um envelope fim-a-fim
formalmente especificado (AES-256-GCM, GMAC ou Ed25519, com o tópico e o cabeçalho
nos dados associados, frescura, janela anti-repetição, política de modos e
restrição chave--tópico) e limitação de tráfego antes de qualquer operação
criptográfica. A proposta não pretende substituir trabalhos com objetivos distintos,
como o anonimato ou a proteção do \emph{broker} por TEE, mas combinar numa só
camada propriedades que a literatura trata em separado.

\subsection{Trabalho futuro}

\begin{itemize}[leftmargin=*]
  \item \textbf{Avaliação experimental (segunda parte do projeto)}: medir, para cada
    combinação de camadas, o custo (latência, processamento, octetos) e o benefício
    (ataques impedidos) sobre um \emph{broker} real, respondendo à lacuna L5.
  \item \textbf{Verificação formal} das propriedades do envelope, à semelhança
    de~\cite{javadi2026cmap}.
  \item \textbf{Acordo de chaves por dispositivo} (por exemplo, ECDH autenticado),
    eliminando chaves pré-partilhadas, e persistência do estado anti-repetição.
  \item \textbf{Cifras para nós sem AES em \emph{hardware}}, como o
    ChaCha20-Poly1305~\cite{rfc8439, sadio2019lightweight}, e cifras resistentes à
    reutilização de \emph{nonces} para chaves de longa duração.
  \item \textbf{Integração com deteção de intrusões}~\cite{vaccari2020mqttset,
    alotaibi2024secure}, usando as rejeições de cada camada como sinal.
  \item \textbf{Atualização periódica da revisão}, reaplicando o protocolo da
    Secção~\ref{sec:metodologia}.
\end{itemize}
